\section{Method}
\label{sec:method}
\subsection{Framework Overview}
\label{sec:overview}

We propose \textbf{DynaDebate}, a multi-agent reasoning framework designed to break initialization homogeneity and ensure process correctness. Formally, given a query $q$, our objective is to derive the correct final answer $y$. The system consists of a set of agents $\mathcal{A} = \{a_1, a_2, \dots, a_N\}$, which interact through a three-stage workflow as illustrated in Figure \ref{fig:framework} (pseudocode in Appendix~\ref{app:pseudocode}).

\subsection{Dynamic Path Generation and Allocation}
\label{sec:path_gen}

Current frameworks often rely on unguided initialization, rendering agents susceptible to generating identical reasoning paths, making it difficult to identify and correct shared errors. To mitigate this, DynaDebate introduces a structured initialization phase comprising Path Generation and Adaptive Redundancy Allocation.

\textbf{Path Generation.}
Prior to the debate, the system utilizes a Path Generation Agent, denoted as $\Phi_{\text{gen}}$, to perform a problem-specific analysis. Given the query $q$, the agent generates a set of candidate reasoning paths $\mathcal{P}$:
\begin{equation}
\mathcal{P} = \Phi_{\text{gen}}(q) = \{p_1, p_2, \ldots, p_K\},
\end{equation}
where $1 \le K \le N$.

The generation process is strictly guided by two constraints:
\begin{enumerate}
    \item \textbf{Logical Soundness:} Each candidate path $p_k$ must be a logically sound method specifically tailored to solving $q$. Logically unsound paths would introduce additional noise into the debate process.
    \item \textbf{Mutual Independence:} To avoid high similarity in agent reasoning and ensure system heterogeneity, we require reasoning paths to be mutually independent. Specifically, for computational tasks, agents employ distinct mathematical methodologies; for information retrieval tasks, they access diverse knowledge sources; and for logical reasoning tasks, they construct unique evidential chains. This constraint forces agents to deviate from the standard response patterns inherent in their training data, thereby ensuring the initial diversity of the solution space.
\end{enumerate}

\textbf{Adaptive Redundancy Allocation.}
To handle cases where the number of valid paths $K$ varies, we employ a Round-Robin Assignment mechanism. The path assigned to agent $a_i$ is defined as:
\begin{equation}
    \text{Path}(a_i) = p_{(i-1 \pmod K) + 1.}
\end{equation}
This formulation implements adaptive redundancy, balancing exploration diversity with consistency verification.
Specifically, in Exploration Mode ($K \approx N$), when multiple distinct paths are feasible, agents are assigned unique paths to encourage the exploration of diverse solutions.
In contrast, under Consistency Check Mode ($K < N$), where valid paths are limited, the same path is assigned to multiple agents, leveraging LLM stochasticity to detect and filter sporadic errors.

\subsection{Process-Centric Debate}
\label{sec:debate}

Following the initialization phase, agents engage in a multi-round debate. Unlike existing debate frameworks that primarily focus on final answers to reach a consensus, DynaDebate enforces a process-centric protocol. This protocol mandates that agents explicitly detail their reasoning process, enabling a rigorous step-by-step audit of the logic presented by their peers.

\textbf{Path-Guided Execution.}
In the initial round ($t=1$), each agent $a_i$ generates a solution strictly adhering to its assigned path $\text{Path}(a_i)$. Agents are prompted to decompose the problem-solving into a sequence of atomic inference steps. Formally, let $R_{i,1}$ denote the reasoning chain generated by agent $a_i$. We define $R_{i,1}$ not as unstructured text, but as an ordered set of logical steps:
\begin{equation}
    R_{i,1} = \{z_i^{(1)}, z_i^{(2)}, \dots, z_i^{(L)}\} \sim \pi( \cdot \mid q, \text{Path}(a_i)),
\end{equation}
where each $z_i^{(k)}$ represents a distinct calculation or logical deduction derived from the preceding steps $z_i^{(<k)}$. This explicit decomposition ensures that the reasoning trajectory is transparent, laying the foundation for the subsequent verification phase.

\textbf{Step-Level Peer Review.}
In subsequent rounds ($t>1$), agents receive the solutions generated by their peers. The core innovation of DynaDebate is the shift in the verification target. Debating agents no longer place their focus on the final answers generated by peers, nor do they assess the correctness of the reasoning process based on structural completeness or textual fluency. Instead, they are prompted to conduct a \textbf{First-Principles Audit} on the inference steps.

Specifically, for a peer's reasoning chain $R_j$, agent $a_i$ performs a rigorous examination targeting both \textbf{intrinsic correctness} and \textbf{logical coherence}. This involves verifying whether each individual step $z_j^{(k)}$ contains calculation or factual errors, while simultaneously auditing the validity of the logical transition from step $z_j^{(k)}$ to $z_j^{(k+1)}$. Agent $a_i$ identifies the specific location where either a factual mistake or a derivation gap occurs, focusing its critique strictly on these identified flaws. This granular scrutiny forces agents to abandon superficial agreement and instead converge solely on logically sound derivations.

\subsection{Trigger-Based Verification }
\label{sec:verification}

Although process-centric peer review strengthens the scrutiny of peer responses, agents may still remain susceptible to blind conformity. Specifically, even when an agent holds the correct answer, it can be influenced by erroneous peer reasoning and may ultimately abandon its original stance to align with the incorrect consensus \cite{wynn2025talk}. To address this limitation, drawing inspiration from the verification paradigm in VerifiAgent~\cite{han2025verifiagent}, DynaDebate integrates a dedicated Trigger-Based Verification Agent, denoted as $\Phi_{\text{ver}}$. The activation of this agent is conditional, governed by a function $\text{Trigger}(H_t)$: at the end of each debate round, if the participating agents produce divergent final answers, $\Phi_{\text{ver}}$ is activated; conversely, if all agents reach unanimous agreement, the debate terminates early without proceeding to the next round, thereby reducing unnecessary token consumption.

Upon activation, the agent autonomously selects and invokes appropriate tools (e.g., a Python interpreter or a web search engine) to execute the command based on the problem context. Given the history $H_t$, the agent generates a verification command and executes it in an external environment to obtain a deterministic observation $o_t \leftarrow \Phi_{\text{ver}}(q, H_t)$. The Verification Agent integrates this result to formulate a reference opinion, updating the debate history:
\begin{equation}
    H_t \leftarrow H_t \cup \{o_t\}.
\end{equation}
This input serves as a critical reference for the subsequent debate round, enabling participating agents to objectively scrutinize the discrepancies between their own stances and the provided evidence.
